\subsection*{用一根箭头保存一块面的信息}
两个三维向量围成一个平行四边形。要同时保存它的面积与朝向，单独一个数不够，而法向量能提供方向。我们约定输出长度等于面积，方向符合选定的右手规则，就得到叉积。它把一块有向面压缩成一根箭头。

沿 $u$ 再沿 $v$ 绕边界，与倒过来绕，围出同一块面，但绕行方向相反。所以交换输入次序，面积大小不变，法向量必须反向。两输入平行时，面塌缩，叉积为零。这两个极端情形可以用来检查任何候选公式。

\subsection*{让坐标公式接受几何检查}
计算前先列要求：输出与两个输入都垂直；任意一边扩大两倍，输出也扩大两倍。坐标公式中的二阶行列式满足这些要求。与 $u$ 点乘时，各项成对抵消，解释了垂直性，而不只是留下一个需要背的三分量式子。

取 $u=(1,2,0),v=(0,1,3)$，输出 $(6,-3,1)$。分别点乘检查与两边正交，再按右手方向检查符号。正交本身不能区别正反，也不能固定长度，因此几个检查必须合起来使用。

下一节会从体积测量反推这根法向量，把点积、行列式和叉积联系起来。那时也会看清这里依赖三维、内积和定向。抽象向量空间中有加法与数乘，并不意味着自动具有同一个叉积运算。
